\documentclass[aspectratio=169]{beamer}
\usepackage{iftex}
\ifPDFTeX
  \usepackage[T1]{fontenc}
  \usepackage{lmodern}
\fi
\usepackage[czech]{babel}
\usepackage{amsmath}
\usepackage{graphicx}
\usepackage{xcolor}
\usepackage{tikz}
\usetikzlibrary{arrows.meta,positioning}
\newcommand{\CVUTTemplatePath}{../../../utils/presentation-template}
\usepackage{\CVUTTemplatePath/beamerthemeCVUT}
\setbeamertemplate{footline}[frame number]
\beamertemplatenavigationsymbolsempty

% Sdílená šablona používá standardní Latin Modern Sans a nevyžaduje externí fonty.
\newcommand{\LANG}{CZ}
\newcommand{\FRONTPAGE}{dark}

\title{Kompilace kódu a verzování}
\subtitle{Základy programování\\Matouš Cejnek\\U12110, FS, ČVUT v Praze}
\author{}
\institute{}
\date{}
\renewcommand{\headlineA}{}
\renewcommand{\headlineB}{}
\renewcommand{\headlineC}{}

\definecolor{CodeBlue}{RGB}{38,83,125}
\definecolor{SoftGray}{RGB}{242,244,247}
\setbeamerfont{normal text}{size=\large}
\setbeamerfont{frametitle}{size=\LARGE}
\setbeamercolor{block title}{fg=white,bg=CodeBlue}
\setbeamercolor{block body}{fg=black,bg=SoftGray}

\newcommand{\term}[1]{\textcolor{CodeBlue}{\textbf{#1}}}
\newcommand{\flowbox}[1]{\node[draw=CodeBlue,very thick,rounded corners,fill=SoftGray,
  minimum height=1.1cm,minimum width=2.7cm,align=center] {#1};}

\begin{document}

% 0:00--0:01
\begin{frame}
  \clearpage\thispagestyle{empty}\titlepage
\end{frame}

% 0:01--0:03
\begin{frame}{Co si dnes odnesete?}
  \begin{itemize}
    \item Co je \term{program}, \term{algoritmus} a \term{programovací jazyk}?
    \item Jak se ze zdrojového kódu stane běžící program?
    \item Jak se liší \term{kompilace} a \term{interpretace}?
    \item K čemu slouží \term{IDE} a online běhové prostředí?
    \item Proč verzujeme a jak Git přemýšlí o historii projektu?
  \end{itemize}
\end{frame}

% 0:03--0:05
\begin{frame}{Co je algoritmus?}
  \begin{block}{Algoritmus}
    Konečný a jednoznačný postup, který převádí vstup na požadovaný výstup.
  \end{block}
  \begin{itemize}
    \item Popisuje \textbf{řešení problému}, ne konkrétní zápis v jazyce.
    \item Musí mít srozumitelné kroky a po konečném počtu kroků skončit.
    \item Tentýž algoritmus lze zapsat slovně, diagramem nebo programem.
  \end{itemize}
\end{frame}

% 0:05--0:07
\begin{frame}{Co je program?}
  \begin{block}{Program}
    Zápis algoritmu a souvisejících dat v podobě, kterou lze vykonat počítačem.
  \end{block}
  \begin{itemize}
    \item \term{Programovací jazyk} určuje povolenou syntax a význam konstrukcí.
    \item \term{Zdrojový kód} je text vytvořený programátorem.
    \item Procesor přímo vykonává až \term{strojové instrukce}.
  \end{itemize}
\end{frame}

% 0:07--0:09
\begin{frame}{Proč počítač nespustí zdrojový kód přímo?}
  \centering
  \begin{tikzpicture}[node distance=1.2cm,>=Latex]
    \node[draw=CodeBlue,very thick,rounded corners,fill=SoftGray,minimum height=1.1cm,
      minimum width=2.7cm,align=center] (source) {Zdrojový\\kód};
    \node[draw=CodeBlue,very thick,rounded corners,fill=SoftGray,minimum height=1.1cm,
      minimum width=3.2cm,right=of source,align=center] (translator) {Překlad nebo\\interpretace};
    \node[draw=CodeBlue,very thick,rounded corners,fill=SoftGray,minimum height=1.1cm,
      minimum width=2.7cm,right=of translator,align=center] (cpu) {Instrukce\\procesoru};
    \draw[->,very thick] (source) -- (translator);
    \draw[->,very thick] (translator) -- (cpu);
  \end{tikzpicture}
  \vspace{0.7cm}
  \begin{itemize}
    \item Zdrojový jazyk je navržen především pro lidi.
    \item Procesor rozumí instrukční sadě konkrétní architektury.
    \item Mezi nimi musí existovat překladač nebo běhové prostředí.
  \end{itemize}
\end{frame}

% 0:09--0:12
\begin{frame}{Co je kompilace?}
  \begin{block}{Kompilace}
    Překlad zdrojového kódu do jiné reprezentace před jeho spuštěním.
  \end{block}
  \begin{itemize}
    \item Kompilátor odhaluje část chyb ještě před spuštěním.
    \item Překlad může probíhat v několika navazujících krocích.
    \item Výsledný program může být závislý na operačním systému a procesoru.
  \end{itemize}
\end{frame}

\begin{frame}{Jak vypadá překladový řetězec?}
  \centering
  \begin{tikzpicture}[node distance=0.65cm,>=Latex]
    \node[draw,rounded corners,fill=SoftGray] (src) {zdrojový kód};
    \node[draw,rounded corners,fill=SoftGray,right=of src] (compiler) {kompilátor};
    \node[draw,rounded corners,fill=SoftGray,right=of compiler] (obj) {objektový kód};
    \node[draw,rounded corners,fill=SoftGray,right=of obj] (linker) {linker};
    \node[draw,rounded corners,fill=SoftGray,right=of linker] (exe) {spustitelný soubor};
    \draw[->] (src)--(compiler); \draw[->] (compiler)--(obj);
    \draw[->] (obj)--(linker); \draw[->] (linker)--(exe);
  \end{tikzpicture}
  \vspace{0.8cm}
  \begin{itemize}
    \item Každý nástroj řeší jinou část převodu programu.
    \item Konkrétní řetězec závisí na jazyce, překladači a platformě.
  \end{itemize}
\end{frame}

\begin{frame}{Co vytváří kompilátor?}
  \begin{block}{Objektový kód}
    Přeložená podoba jedné části programu, která ještě nemusí být samostatně spustitelná.
  \end{block}
  \begin{itemize}
    \item Obsahuje strojové instrukce a data pro cílovou platformu.
    \item Může stále odkazovat na funkce nebo data definovaná jinde.
    \item Větší projekt se obvykle překládá po samostatných částech.
  \end{itemize}
\end{frame}

\begin{frame}{Co dělá linker?}
  \begin{block}{Linker neboli sestavovací program}
    Spojuje objektové soubory a knihovny do výsledného programu.
  \end{block}
  \begin{itemize}
    \item Vyhledá, kde jsou definované používané funkce a proměnné.
    \item Propojí odkazy mezi jednotlivými částmi programu.
    \item Ohlásí chybu, pokud potřebná definice chybí nebo je nejednoznačná.
  \end{itemize}
\end{frame}

\begin{frame}{Co je knihovna?}
  \begin{block}{Knihovna}
    Znovupoužitelný soubor již připraveného kódu, který program využívá.
  \end{block}
  \begin{itemize}
    \item \term{Statická knihovna}: potřebné části se vloží do programu při linkování.
    \item \term{Dynamická knihovna}: načte se při spuštění nebo až během běhu.
    \item Knihovna není samostatná aplikace; poskytuje funkce jiným programům.
  \end{itemize}
\end{frame}

\begin{frame}{Co se děje při spuštění?}
  \begin{itemize}
    \item Operační systém načte spustitelný soubor do paměti.
    \item \term{Loader} připraví program a potřebné dynamické knihovny.
    \item Runtime inicializuje prostředí, které program potřebuje.
    \item Řízení přejde na vstupní bod programu.
  \end{itemize}
  \begin{alertblock}{Spustitelný neznamená přenositelný}
    Soubor vytvořený pro jednu platformu nemusí fungovat na jiné.
  \end{alertblock}
\end{frame}

% 0:12--0:15
\begin{frame}{Co je interpretace?}
  \begin{block}{Interpretace}
    Interpret načítá program a zajišťuje vykonávání jeho instrukcí za běhu.
  \end{block}
  \begin{itemize}
    \item Ke spuštění potřebujeme zdrojový kód a vhodný interpret/runtime.
    \item Chyba se může projevit až při vykonání problematické části programu.
    \item Typický příklad uživatelského pohledu: \texttt{python program.py}.
  \end{itemize}
  \begin{alertblock}{Pozor}
    „Interpretovaný jazyk“ neznamená, že uvnitř neprobíhá žádný překlad.
  \end{alertblock}
\end{frame}

\begin{frame}{Interpret a virtuální stroj}
  \begin{itemize}
    \item \term{Interpret} zajišťuje vykonávání programu za běhu.
    \item \term{Virtuální stroj} vykonává instrukce mezireprezentace, například bytecode.
    \item Stejná mezireprezentace může běžet na více platformách s vhodným runtime.
  \end{itemize}
  \begin{block}{Názvy popisují vrstvy, ne vždy celý jazyk}
    Jeden jazyk může mít více implementací s odlišným způsobem vykonávání.
  \end{block}
\end{frame}

% 0:15--0:18
\begin{frame}{Kompilace, nebo interpretace?}
  \begin{columns}[T]
    \begin{column}{0.47\textwidth}
      \textbf{Překlad před spuštěním}
      \begin{itemize}
        \item vzniká cílový kód,
        \item více práce předem,
        \item často rychlý běh.
      \end{itemize}
    \end{column}
    \begin{column}{0.47\textwidth}
      \textbf{Překlad během běhu}
      \begin{itemize}
        \item nutné běhové prostředí,
        \item snadná přenositelnost,
        \item pružná práce s programem.
      \end{itemize}
    \end{column}
  \end{columns}
  \vspace{0.5cm}
  \begin{block}{V praxi jde o spektrum}
    Jazyky a implementace mohou kombinovat kompilaci, bytecode, virtuální stroj i JIT překlad.
  \end{block}
\end{frame}

% 0:18--0:20
\begin{frame}{Kde je v tomto obrazu Python?}
  \begin{itemize}
    \item Python je \term{programovací jazyk}; CPython je jeho běžná implementace.
    \item CPython typicky překládá zdrojový kód do \term{bytecode}.
    \item Bytecode vykonává virtuální stroj uvnitř Python runtime.
  \end{itemize}
  \centering
  \vspace{0.5cm}
  \begin{tikzpicture}[node distance=0.9cm,>=Latex]
    \node[draw,rounded corners,fill=SoftGray] (py) {\texttt{program.py}};
    \node[draw,rounded corners,fill=SoftGray,right=of py] (bc) {bytecode};
    \node[draw,rounded corners,fill=SoftGray,right=of bc] (vm) {Python VM};
    \node[draw,rounded corners,fill=SoftGray,right=of vm] (out) {výsledek};
    \draw[->] (py)--(bc); \draw[->] (bc)--(vm); \draw[->] (vm)--(out);
  \end{tikzpicture}
\end{frame}

% 0:20--0:22
\begin{frame}{Co je runtime?}
  \begin{block}{Běhové prostředí (runtime)}
    Služby a pravidla potřebná při spuštění programu.
  \end{block}
  \begin{itemize}
    \item načtení programu a knihoven,
    \item správa paměti a objektů,
    \item vstup, výstup a komunikace s operačním systémem,
    \item zpracování běhových chyb.
  \end{itemize}
  \vspace{0.3cm}
  \textbf{Jazyk} je specifikace; \textbf{implementace} a runtime ji uskutečňují.
\end{frame}

% 0:22--0:25
\begin{frame}{Chyby před spuštěním}
  \begin{itemize}
    \item \term{Syntaktická chyba}: zápis neodpovídá pravidlům jazyka.
    \item \term{Překladová chyba}: program nelze převést do cílové podoby.
    \item \term{Linkovací chyba}: některou požadovanou definici nelze najít nebo propojit.
  \end{itemize}
  \begin{block}{Výhoda včasné kontroly}
    Část problémů lze najít ještě dříve, než program dostane vstup a začne pracovat.
  \end{block}
\end{frame}

\begin{frame}{Chyby během a po spuštění}
  \begin{itemize}
    \item \term{Běhová chyba}: problém vznikne až během vykonávání.
    \item \term{Logická chyba}: program běží, ale dává nesprávný výsledek.
  \end{itemize}
  \begin{alertblock}{Důležitý rozdíl}
    Úspěšný překlad ani spuštění nedokazují, že program řeší správný problém.
  \end{alertblock}
\end{frame}

% 0:25--0:27
\begin{frame}{Co je IDE?}
  \begin{block}{IDE --- integrované vývojové prostředí}
    Aplikace, která spojuje nástroje pro tvorbu a zkoumání programu.
  \end{block}
  \begin{itemize}
    \item editor kódu a navigace v projektu,
    \item zvýraznění syntaxe a průběžná analýza,
    \item spuštění, sestavení a práce s runtime,
    \item debugger, testy, terminál a integrace verzování.
  \end{itemize}
\end{frame}

% 0:27--0:29
\begin{frame}{Textový editor, nebo IDE?}
  \begin{itemize}
    \item \term{Textový editor}: především úprava souborů; další nástroje doplňujeme.
    \item \term{IDE}: ucelené prostředí s přednastaveným workflow projektu.
  \end{itemize}
  \begin{block}{Volba nástroje}
    Nástroj mění způsob práce, nikoli význam programu ani pravidla jazyka.
  \end{block}
\end{frame}

\begin{frame}{Notebook a cloudové prostředí}
  \begin{itemize}
    \item \term{Notebook}: dokument kombinující text, kód a jednotlivé výstupy.
    \item \term{Cloudové IDE / Colab}: běh na vzdáleném počítači přes prohlížeč.
    \item Výhodou je rychlý začátek; omezením může být závislost na službě a připojení.
  \end{itemize}
\end{frame}

% 0:29--0:31
\begin{frame}{Proč verzovat?}
  \begin{itemize}
    \item Víme, \textbf{co}, \textbf{kdo} a \textbf{proč} změnil.
    \item Můžeme porovnat stavy projektu a vrátit se ke starší verzi.
    \item Více lidí může bezpečně rozvíjet tentýž projekt.
    \item Experiment lze oddělit od stabilní práce.
  \end{itemize}
  \begin{alertblock}{Složky \texttt{final}, \texttt{final2}, \texttt{opravdu-final}}
    Kopie souborů nejsou spolehlivá historie: chybí vztahy, autor, důvod změny i možnost rozumného sloučení.
  \end{alertblock}
\end{frame}

% 0:31--0:33
\begin{frame}{Co je Git?}
  \begin{block}{Git}
    Distribuovaný systém správy verzí, který uchovává historii projektu jako propojené snímky.
  \end{block}
  \begin{itemize}
    \item \term{Repozitář} obsahuje soubory projektu i jejich historii.
    \item Každý klon má obvykle úplnou lokální historii.
    \item Většina operací funguje lokálně, bez připojení k serveru.
    \item Git není totéž co GitHub: Git je nástroj, GitHub je hostingová služba.
  \end{itemize}
\end{frame}

\begin{frame}{Git ukládá snímky, ne názvy verzí}
  \begin{itemize}
    \item Každý commit představuje snímek stavu sledovaných souborů.
    \item Nezměněný obsah Git nemusí ukládat znovu; odkáže na existující data.
    \item Commity jsou propojené a vytvářejí historii projektu.
  \end{itemize}
  \centering
  \vspace{0.5cm}
  \texttt{A --- B --- C --- D}
\end{frame}

% 0:33--0:36
\begin{frame}{Jak Git vidí změnu?}
  \centering
  \begin{tikzpicture}[node distance=1.6cm,>=Latex]
    \node[draw,rounded corners,fill=SoftGray,minimum width=2.6cm] (work) {pracovní strom};
    \node[draw,rounded corners,fill=SoftGray,minimum width=2.6cm,right=of work] (stage) {staging area};
    \node[draw,rounded corners,fill=SoftGray,minimum width=2.6cm,right=of stage] (repo) {repozitář};
    \draw[->,very thick] (work)--node[above]{\texttt{add}}(stage);
    \draw[->,very thick] (stage)--node[above]{\texttt{commit}}(repo);
  \end{tikzpicture}
  \vspace{0.8cm}
  \begin{itemize}
    \item \term{Modified}: soubor je změněný v pracovním stromu.
    \item \term{Staged}: konkrétní podoba změny je vybraná pro další commit.
    \item \term{Committed}: snímek je bezpečně uložený v lokální historii.
  \end{itemize}
\end{frame}

% 0:36--0:38
\begin{frame}{Co je commit?}
  \begin{block}{Commit}
    Pojmenovaný snímek projektu s autorem, časem, zprávou a odkazem na předchozí historii.
  \end{block}
  \begin{itemize}
    \item Zachycuje jeden smysluplný krok vývoje.
    \item Identifikuje jej hash odvozený z jeho obsahu a metadat.
    \item Dobrá zpráva vysvětluje záměr změny, nejen jména souborů.
  \end{itemize}
  \vspace{0.3cm}
  \centering
  \texttt{A --- B --- C} \hspace{0.4cm} \textcolor{CodeBlue}{$\leftarrow$ historie snímků}
\end{frame}

% 0:38--0:40
\begin{frame}{Co je větev?}
  \vspace{0.35cm}
  \term{Větev (branch)} je pohyblivý ukazatel na commit; umožňuje rozvíjet více linií historie.\par
  \centering
  \begin{tikzpicture}[>=Latex]
    \node[circle,draw,fill=SoftGray] (a) at (0,0) {A};
    \node[circle,draw,fill=SoftGray] (b) at (1.7,0) {B};
    \node[circle,draw,fill=SoftGray] (c) at (3.4,0) {C};
    \node[circle,draw,fill=SoftGray] (d) at (5.1,0) {D};
    \node[circle,draw,fill=SoftGray] (x) at (3.4,-1.0) {X};
    \draw[->] (a)--(b); \draw[->] (b)--(c); \draw[->] (c)--(d);
    \draw[->] (b)--(x);
    \node[above=0.15cm of d] {\texttt{master}};
    \node[below=0.15cm of x] {\texttt{experiment}};
  \end{tikzpicture}
  \begin{itemize}
    \item Vytvoření větve nekopíruje celý projekt.
    \item \term{Merge} znovu propojí dvě linie vývoje.
  \end{itemize}
\end{frame}

\begin{frame}{Co je merge a konflikt?}
  \begin{itemize}
    \item \term{Merge} spojí změny ze dvou linií historie.
    \item Nezávislé změny Git často sloučí automaticky.
    \item \term{Konflikt} vznikne, když nelze bezpečně rozhodnout, kterou podobu zachovat.
    \item Konflikt není porucha Gitu; je to rozhodnutí, které musí udělat člověk.
  \end{itemize}
\end{frame}

% 0:40--0:42
\begin{frame}{Lokální a vzdálený repozitář}
  \begin{itemize}
    \item \term{Lokální repozitář}: pracovní kopie a historie na našem počítači.
    \item \term{Vzdálený repozitář}: další kopie historie, často na GitHubu.
    \item \texttt{clone}: vytvoří lokální kopii vzdáleného repozitáře.
    \item \texttt{push}: zveřejní lokální commity; \texttt{pull}: převezme vzdálené změny.
  \end{itemize}
  \begin{alertblock}{Častý omyl}
    Commit ukládá změnu lokálně. Teprve push ji odešle na vzdálený server.
  \end{alertblock}
\end{frame}

\begin{frame}{Fork není clone}
  \begin{columns}[T]
    \begin{column}{0.46\textwidth}
      \begin{block}{Fork}
        Kopie cizího repozitáře pod naším účtem na GitHubu.
      \end{block}
      \begin{itemize}
        \item vzniká na serveru,
        \item můžeme do ní zapisovat,
        \item zůstává propojená s původním projektem.
      \end{itemize}
    \end{column}
    \begin{column}{0.46\textwidth}
      \begin{block}{Clone}
        Lokální kopie vzdáleného repozitáře v našem počítači.
      \end{block}
      \begin{itemize}
        \item vzniká příkazem \texttt{git clone},
        \item obsahuje soubory i historii,
        \item umožňuje lokálně upravovat a commitovat.
      \end{itemize}
    \end{column}
  \end{columns}
  \begin{alertblock}{Pořadí pro cvičení}
    Nejdřív vytvoříme fork na GitHubu, potom naklonujeme svůj fork.
  \end{alertblock}
\end{frame}

\begin{frame}{Proč potřebujeme origin a upstream?}
  \centering
  \begin{tikzpicture}[>=Latex]
    \node[draw,rounded corners,fill=SoftGray,minimum width=2.7cm,align=center]
      (upstream) at (0,0) {repozitář vyučujícího\\\texttt{upstream}};
    \node[draw,rounded corners,fill=SoftGray,minimum width=2.7cm,align=center]
      (local) at (4,0) {lokální repozitář\\v počítači};
    \node[draw,rounded corners,fill=SoftGray,minimum width=2.7cm,align=center]
      (origin) at (8,0) {náš fork\\\texttt{origin}};
    \draw[->,very thick] (upstream)--(local);
    \draw[<->,very thick] (local)--(origin);
  \end{tikzpicture}
  \vspace{0.4cm}
  \begin{itemize}
    \item \term{Remote} je pojmenovaný odkaz na vzdálený repozitář.
    \item \texttt{origin} ukazuje na náš fork; \texttt{upstream} na původní repozitář.
    \item Z \texttt{upstream} změny načítáme; do \texttt{origin} je odesíláme.
    \item Adresy vždy ověříme pomocí \texttt{git remote -v}.
  \end{itemize}
\end{frame}

\begin{frame}{Jak synchronizovat master?}
  \begin{enumerate}
    \item Přepneme se na lokální hlavní větev: \texttt{git switch master}.
    \item Načteme stav původního projektu: \texttt{git fetch upstream}.
    \item Začleníme změny: \texttt{git merge upstream/master}.
    \item Aktualizujeme svůj fork: \texttt{git push origin master}.
  \end{enumerate}
  \texttt{fetch} pouze načte vzdálenou historii; teprve \texttt{merge} změní
  naši aktuální větev.
  \vspace{0.2cm}
  \begin{alertblock}{Před založením pracovní větve}
    Novou větev vytváříme z aktualizovaného \texttt{master}, jinak začínáme
    pracovat na zastaralém základu.
  \end{alertblock}
\end{frame}

\begin{frame}{Co je pull request?}
  \begin{block}{Pull request}
    Návrh na začlenění jedné větve do druhé.
  \end{block}
  \begin{itemize}
    \item Ukazuje rozdíl mezi zdrojovou a cílovou větví.
    \item Umožňuje změny před sloučením zkontrolovat a okomentovat.
    \item Na GitHubu může spojovat větve ve stejném i v různých repozitářích.
  \end{itemize}
  \begin{alertblock}{Pozor na cílový repozitář}
    V tomto cvičení cílíme na svůj fork, ne na repozitář vyučujícího.
  \end{alertblock}
\end{frame}

\begin{frame}{Kam poslat změny?}
  \begin{enumerate}
    \item Z aktualizovaného \texttt{master} vytvoříme pracovní větev.
    \item Změny commitujeme lokálně do pracovní větve.
    \item Větev odešleme pomocí \texttt{git push} do svého \texttt{origin}.
    \item Na GitHubu vytvoříme pull request do \texttt{master} svého forku.
  \end{enumerate}
  \vspace{0.35cm}
  \textbf{Před potvrzením:} Zkontrolujeme zdrojovou i cílovou větev a oba
  repozitáře.
\end{frame}

% 0:42--0:44
\begin{frame}{Jeden projekt, dvě vrstvy nástrojů}
  \begin{columns}[T]
    \begin{column}{0.47\textwidth}
      \textbf{Vznik a běh programu}
      \begin{itemize}
        \item zdrojový kód,
        \item překladač / interpret,
        \item runtime,
        \item IDE.
      \end{itemize}
    \end{column}
    \begin{column}{0.47\textwidth}
      \textbf{Historie projektu}
      \begin{itemize}
        \item pracovní strom,
        \item staging area,
        \item commity a větve,
        \item vzdálený repozitář.
      \end{itemize}
    \end{column}
  \end{columns}
  \vspace{0.5cm}
  IDE může obě vrstvy zpřístupnit, ale nenahrazuje jejich pochopení.
\end{frame}

% 0:44--0:45
\begin{frame}{Co si zapamatovat?}
  \begin{itemize}
    \item Zdrojový kód musí být přeložen nebo vykonán prostřednictvím runtime.
    \item Kompilace a interpretace nejsou dvě dokonale oddělené kategorie.
    \item IDE spojuje nástroje, ale programovací jazyk ani Git se na něj nevážou.
    \item Git ukládá propojené snímky projektu v commitech.
    \item Fork je naše serverová kopie; clone je lokální kopie v počítači.
    \item \texttt{origin} míří na náš fork, \texttt{upstream} na původní projekt.
  \end{itemize}
  \vspace{0.4cm}
  \centering\Large Otázky?
\end{frame}

\end{document}
